\documentclass[english,twoside, openright, hidelinks, 11 pt, class=report,crop=false]{standalone}
\input{../../preamb}
\input{../bm}

\begin{document}

\nes

\op{ggb1}
\textbf{a)} Skriv den lineære funksjonen $ {f(x)=2x+4} $ og linja $ {y=2x+2} $ inn i GeoGebra. Lag $ f(x) $ blå og $ y $ grønn. Hva ser du ut ifra grafen til de to linjene?\bs
\textbf{b)} Finn verdien til $ f(x) $ når $ {x=4} $.\bs
\textbf{c)} Finn verdien til $ y $ når $ {x=-3} $.\vsk

\op{ggb2}
\textbf{a)} Tegn punktene $ (-1,2) $ og $ (2,8) $.\bs
\textbf{b)} Finn uttrykket til linja som går gjennom disse punktene.\vsk

\op{ggb3}
\textbf{a)} Skriv inn funksjonen $ {f(x)=x^2+2x-3} $.\bs
\textbf{b)} Finn $ f(4) $.\bs
\textbf{c)} Finn nullpunktene til $ f(x) $.\bs
\textbf{d)} Finn bunnpunktet til $ f(x) $.\bs
\textbf{e)} Finn skjæringspunktet mellom $ f(x) $ og linja $ y=5 $.

\eksop{gv25d2opg6}{GV25D2}
Elevrådet skal arrangere avslutningsfest for 10.trinn og undersøker pris for mat og drikke og leie av lokale. De har funnet tre tilbud: \vs
\begin{center}
	\begin{tabular}{c | l}
		\hline
		\textbf{Tilbud 1} & 150 kroner per elev for mat, drikke og leie av lokale. \\ \hline
		\textbf{Tilbud 2} & 75 kroner per elev for mat og drikke, \\
		& pluss 3000 kroner i leie av lokale.  \\ \hline
		\textbf{Tilbud 3} & 7\,000 kroner som dekker mat, drikke og lokale \\
		&for inntil 70 elever. \\ \hline
	\end{tabular}
\end{center}
\abc{
\item Lag en grafisk fremstilling av de tre tilbudene i samme koordinatsystem.
\item Sammenlign tilbudene, og finn ut når elevrådet bør velge de ulike tilbudene.
}
\newpage

\grubeksoprecc{gv24d1opg2}{GV24D2}
Gitt funksjonene
\[ f(x)=200x+4 \quad,\quad, g(x)=\frac{1000}{x}+80\quad,\quad-100x\]
\abc{
\item Tegn grafene til funksjonene, og forklar hvilke typer funskoner de er, og hva som kjennetegner dem.
\item Velg én av funksjonene, og gi eksempel på en praktisk\\ situasjon som funksjonen kunne ha beskrevet.
}

\grubop{opgstatscriptparodd}
Lag et skript som fra en liste med tall finner
\abc{
	\item gjennomsnittet.
	\item medianen.
	\item typetallet.	
}
Undersøk svarene fra {\color{blue}oppgave \ref{opgstatsentrodd1} - \ref{opgstatsentrpar2}} med sriptet.


\end{document}

